
多项式根式可解(Solvable by Radicals)
已知$f(x)$的分裂域为$E/F$, 称其为根式可解
如果存在根塔 $F=F_1\subseteq\cdots\subseteq F_{r+1}=K$
满足 $F_{i+1}=F_i(d_i)\land d_i^{n_i}\in F_i\land E\subseteq K$


------

引理1: $K/F\implies G_K(f)<G_F(f)$
证明: 已知$f(x)$的根集为$\{a_1,\cdots,a_n\}$
构造$L=K(a_1,\cdots,a_n)$, 并且$L\supseteq E$
可构造群同态$\phi:\text{Gal}(L/K)\to\text{Gal}(E/F):\eta\mapsto\eta|E$
$\ker(\phi)=\{\eta(a_i)=a_i\land\text{K-自同构}\}=\{e\}$, 故是单同态

------

正规闭包的定义(Normal Closure)
已知代数扩张$E/F$, 则将包含$E$的最小正规扩张, 称为正规闭包

------

引理2: $E/F$存在根塔$\implies$正规闭包$K/F$也存在根塔
证明: 构造$K'=\langle\eta(E)\mid\eta\in\text{Gal}(K/F)\rangle$
那么$\eta(K')\subseteq K'\implies K'/F$也是正规扩张$\implies K=K'=\langle\eta(E)\rangle$
现设$\text{Gal}(K/F)=\{\eta_1,\cdots,\eta_k\}$, 将$\eta_j$作用于原根塔得

$F=\eta_jF_1\subseteq\eta_jF_2\cdots\subseteq\eta_jF_{r+1}=\eta_j(E)$
满足 $\eta_jF_{i+1}=\eta_j[F_i(d_i)]=(\eta_jF_i)(\eta_j d_i)$
其中 $(\eta_jd_i)^{n_i}=\eta_j(d_i^{n_i})\in\eta_j(F_i)$

$\begin{aligned}
    F&\subseteq F(\eta_1d_1)\subseteq F(\eta_1d_1,\eta_1d_2)\subseteq\cdots\subseteq F(\eta_1d_1,\cdots,\eta_1d_r)=F(\eta_1E) \\
     &\subseteq F(\eta_1E,\eta_2d_1)\subseteq\cdots\subseteq F(\eta_1E,\eta_2E)\subseteq\cdots\subseteq F(\eta_1E,\cdots,\eta_kE)=K \\
\end{aligned}$

------

伽罗瓦判别定理(Galois Criterion)
在特征为0的域上, 多项式根式可解$\iff$其分裂域的伽罗瓦群为可解群
